\documentclass[11pt]{article}
\usepackage[a4paper,margin=26mm]{geometry}
\usepackage{amsmath,amssymb,amsthm,mathtools}
\usepackage[T1]{fontenc}
\usepackage{lmodern,microtype}
\usepackage[colorlinks=true,linkcolor=blue,urlcolor=blue,citecolor=blue]{hyperref}
\hypersetup{pdftitle={The Unique Games Theorem: A Short Self-Contained Proof},pdfauthor={Ariel Elboim}}
\usepackage{enumitem}
\setlist{itemsep=3pt,topsep=5pt}
\newtheorem{theorem}{Theorem}
\newtheorem{proposition}[theorem]{Proposition}
\newtheorem{lemma}[theorem]{Lemma}
\newcommand{\E}{\mathbb E}
\newcommand{\F}{\mathbb F}
\newcommand{\eps}{\varepsilon}
\newcommand{\rank}{\operatorname{rank}}
\newcommand{\Hom}{\operatorname{Hom}}
\newcommand{\Tr}{\operatorname{Tr}}
\newcommand{\Span}{\operatorname{span}}
\newcommand{\val}{\operatorname{val}}
\newcommand{\OPT}{\operatorname{OPT}}
\newcommand{\ind}{\mathbf 1}
\newcommand{\depositdoi}{10.5281/zenodo.23265572}
\title{The Unique Games Theorem:\\ A Short Self-Contained Proof}
\author{Ariel Elboim\\\small Independent researcher, ORCID 0009-0001-1341-366X\\\small Prepared with Codex (OpenAI), Claude (Anthropic) and other AI models; see the companion note}
\date{9 October 2026}
\begin{document}
\maketitle
\begin{abstract}
We give a short proof of the theorem of OpenAI's preprint \emph{The
Unique Games Theorem} \cite{UGT}. It is self-contained apart from three
earlier theorems, which are stated as inputs. As in the source, the
argument constructs stable nonlinear noise, extracts constant-output
certificates, decodes them using only local information, and contradicts
parallel repetition on clean coordinates. Relative to the source,
several steps are simplified: the noise gadget's recursion is written in
direct coordinates, and the survival of linear rank is proved by one
induction on the height; genericity is proved by formal differentiation;
test acceptance is bounded directly by the measure of good row fibers,
without recoloring or concentration; the second player guesses a
constant output and Fourier-samples one row; and a triangle of three
equations replaces variable cloning. A companion note \cite{Notes} lists
all departures from the source and their credits. A longer presentation
that follows the source closely is \cite{Elboim}. The proof is written
for a reader comfortable with finite-field linear algebra and Fourier
analysis on finite groups.
\end{abstract}

\noindent\textbf{Status.}
This is a simplified reconstruction of the proof in \cite{UGT}, prepared
with AI agents and checked step by step by separate AI reviewers. These
checks do not certify the preprint or verify this paper, and no human
expert has reviewed it. OpenAI's repository also contains a Lean~4
formalization of the source theorem, whose Comparator check passed when
we ran it; that check concerns OpenAI's formal development, not this
paper or its simplifications. Details of the check, and the history of
this paper, are in the companion note \cite{Notes}. This paper is not
derived from the formalization. Fortnow's earlier exposition
\cite{Fortnow} presents the same proof at greater length; this paper was
prepared independently of it. The theorem below is the preprint's claim,
whose proof we present.

\section{The claim and three prior inputs}

A unique game labels vertices from a finite alphabet. Each edge carries
a permutation: a label at one endpoint specifies exactly one acceptable
label at the other. Its value is the largest fraction of constraints
satisfied by one labeling.

\begin{theorem}[Main claim of \cite{UGT}]\label{thm:main}
For every fixed $0<\eps,\delta<1/2$, there is a deterministic polynomial-time
reduction from satisfiability to explicit unweighted unique games on
simple bipartite graphs that maps satisfiable formulas to YES games and
unsatisfiable formulas to NO games, where
\[
 \mathrm{YES}:\val(G)\ge1-\eps,
 \qquad \mathrm{NO}:\val(G)\le\delta.
\]
The alphabet is a fixed binary vector space, and every constraint is a
translation. The alphabet and the degree of the running-time polynomial
may depend on the two target errors, but not on the input length.
\end{theorem}

All spaces and linear maps below are over $\F_2$, except where a larger
finite field is specified. Uniform samples include zero. Separately
specified random choices are independent unless stated otherwise.
Intermediate games may have rational weights and repeated constraint
occurrences.

\paragraph{Input 1: the parity gap.}
For every fixed positive rational $\xi$, H\aa stad's theorem
\cite[Theorem 5.4 and its proof]{Hastad} gives a deterministic polynomial-time reduction
from satisfiability to lists of three-bit equations, repetitions allowed,
\[
 x_{v_1}+x_{v_2}+x_{v_3}=b_e,
 \qquad \mathrm{YES}:\OPT\ge1-\xi,
 \qquad \mathrm{NO}:\OPT\le\tfrac12+\xi,
\]
where $\OPT$ is the largest satisfiable fraction. The list has one entry
per random string of that proof's verifier. If the formula with three
literals per clause built in that proof has $M$ clauses, the verifier picks
$u$ independent uniform clause--variable pairs, then a fixed budget of $b$
fair bits, padded with unused bits when the chosen clauses share variables;
$u$, the dyadic noise rate and hence $b$ depend only on $\xi$. So its random
strings form a uniform space of size $(3M)^u2^b$, polynomial in the input,
and the list is the unweighted multiset of one equation per string. This
size is the number $N_0$ of entries in Section~\ref{sec:tables}, and gives
the common denominators $3N_0$ and $N_{\mathcal H}$ of
Section~\ref{sec:finish}.

\paragraph{Input 2: Grassmann expansion.}
The binary Grassmann graph has the $\ell$-dimensional subspaces of an
ambient space as vertices; adjacency means an intersection of dimension
$\ell-1$. The retention of a nonempty set is the probability that a uniform
neighbor of a uniform member stays in the set. For every fixed retention
threshold $\zeta>0$, there are $\alpha>0$ and an integer $r_{\rm s}\ge1$, depending
only on $\zeta$, such that every set with retention at least $\zeta$ has
relative density at least $\alpha$ in an interval
\[
 \{L:A_0\subseteq L\subseteq B_0\},
 \qquad \dim A_0+\operatorname{codim}B_0\le r_{\rm s}.
\]
This holds for every sufficiently large $\ell$, and then every sufficiently
large ambient dimension. This is the expansion theorem of
Khot--Minzer--Safra \cite[Theorem 1.12]{KMS}.

\paragraph{Input 3: projection-game repetition.}
In a projection game, an answer of the first player determines one required
answer of the second. Dinur--Steurer \cite[Corollary 1.2]{DS} give
\begin{equation}\label{eq:repetition}
 \val(G)\le1-g\quad\Longrightarrow\quad
 \val(G^{\otimes n})\le(1-g^2/16)^n.
\end{equation}
It applies to finite weighted games, lets each player use its entire
question tuple, and does not depend on the alphabet. Our games have a
fixed projection for each question pair.

\paragraph{Idea and plan.}
The honest labeling applies one nonlinear map $C$, and the test perturbs
its input by a special noise. The noise rarely changes the output of $C$,
yet every linear family that sees enough of the output notices it with
constant probability. Section~\ref{latent:section} builds such noise by a
race inside a recursive construction. At each level an output change
survives with probability $1-\theta$, so after $s$ levels it survives with
probability about $e^{-\theta s}$. For a generic family of linear
observations of rank $r$, the potential $H_r$ drops by at most $3\theta$ per
level in expectation. Like a coupon collector, the family therefore needs
at least about $H_r/(3\theta)$ levels to die, where
$H_r=1+\frac12+\dots+\frac1r$. The harmonic numbers diverge, so a large
initial rank wins; the proof takes $s\approx H_{r_0}/(8\theta)$. The algebra
needed is that, for a generic family, at most $3r$ elements are aligned.
An aligned count of order $2^r$, as for arbitrary spaces, would make the
waiting times $\sum_j1/(2^j\theta)$ summable, and this argument would fail.
Section~\ref{sec:tables} builds the translation game and proves its
completeness, Section~\ref{sec:decoders} proves its soundness
(Proposition~\ref{prop:sound}), and Section~\ref{sec:finish} removes the
weights and amplifies the gap.

\section{From parity equations to shared tables}\label{sec:tables}

\paragraph{Distinct variables and the base gap.} Let $\mathcal E_0$ be the
list from Input~1, each of its $N_0$ entries of weight $1/N_0$. Over $\F_2$
a repeated variable cancels, so each entry has three distinct variables or
is unary, $x=b$. Replace each unary occurrence of weight $w$ by the three
equations $x+z_1+z_2=b$, $x+z_2+z_3=b$, $x+z_1+z_3=b$ of weight $w/3$, with
fresh $z_1,z_2,z_3$. If $x=b$, equal auxiliary bits satisfy all three; if
$x\ne b$, all three would need $z_i+z_j=1$, impossible since the three
pairwise sums add to $0$ while $1+1+1=1$, so at most two hold. Hence every
occurrence of weight $w$, replaced or not, contributes at most
$w(\frac23+\frac13\ind)$, where $\ind$ indicates that the original
occurrence is satisfied. Write $\mathcal E$ for the processed system with
weights $\omega_e$. Extending by zero auxiliary bits gives $\OPT(\mathcal
E)\ge\OPT(\mathcal E_0)$, and summing gives $\OPT(\mathcal
E)\le\frac23+\frac13\OPT(\mathcal E_0)$, so in the NO case, for
$\xi<1/100$, $\OPT(\mathcal E)\le\frac56+\frac\xi3<6/7$. In the
equation--variable game the first player receives an equation and returns a
satisfying triple; the second receives a uniformly chosen one of its three
variables and returns a bit. Distinctness makes the pair specify a unique
position. The second player's answers form a global assignment. On each equation it
violates, every legal triple disagrees with it in some position, which is
asked with probability $1/3$; so in the NO case
\begin{equation}\label{eq:base-gap}
 \val(G_{\mathcal E})\le1-\frac{1-\OPT(\mathcal E)}3\le20/21.
\end{equation}

\paragraph{The noise gadget needed below.}
Section~\ref{latent:section} constructs, for each $p_*>0$, a threshold $r_*$
such that for every sufficiently large $\ell$ divisible by $3$ there are
finite data
\begin{equation}\label{eq:gadget-contract}
 \begin{gathered}
 K=\F_2^\ell\le\mathcal V,\quad C:\mathcal V\to K,\quad a\sim\mu,\\
 C(x+h)=C(x)+h\quad(h\in K),\qquad
 \Pr_{x\text{ uniform},a}[C(x+a)\ne C(x)]\le p_*,\\
 G\le\mathcal V^*,\ \dim(G|_K)\ge r_*
 \quad\Longrightarrow\quad
 \Pr_a[\exists g\in G:g(a)\ne0]\ge1/8.
 \end{gathered}
\end{equation}
The noise law is rational. Crucially, $r_*$ is chosen before $\ell$.
Fix a linear splitting $\mathcal V=K\oplus W$.

\paragraph{Questions and homogeneous answers.}
Sample $U=(e_1,\ldots,e_k)$ independently with replacement. Put
$L_U=\F_2^{1+2k}$, with coordinates $t_0,x_1,y_1,\ldots,x_k,y_k$.
The triple for equation $e_i$ is
\[
 (x_i,y_i,b_{e_i}t_0+x_i+y_i).
\]
Vectors with $t_0=1$ are exactly the legal answers. Different positions
retain separate answer coordinates, even when their variable identifiers
repeat. A projected question $O$ retains either an equation or one of its
actual bits in each position; a singleton position records only the
variable's identifier, not the equation or the slot it came from. This gives a surjective linear map
$\pi:L_U\to L_O$ preserving $t_0$.

\paragraph{Table keys force sharing.}
A table is a linear map $P:L_Q\to\mathcal V$, or an affine map on the
legal answers. On the $t_0=0$ part of an equation position, the three
answer bits $x_i,y_i$ and $z_i=b_{e_i}t_0+x_i+y_i$ restrict to exactly the
nonzero functionals; a singleton position has only its one bit. So in each position $P$ has rank two, and depends on
the whole equation; or rank one, and is $v$ times one answer bit; or
vanishes. Hence $P$
factors as $P'\pi'$ through the coarsest projected question that keeps,
in each position, exactly what $P$ depends on, omitting the positions where
$P$ vanishes. Its key is this
factorization: positions, types, identifiers and the map $P'$. Pulling a
projected table back to $U$ and reducing again gives the same key: $\pi$
only deletes answer bits, so the pullback depends on the same bits, its
factorization map is that of the projected table composed with $\pi$, and
$P'$ is unique. Answer bits are functionals on $L_U$: at a singleton
position keeping the third bit, $\pi$ keeps $z_i=b_{e_i}t_0+x_i+y_i$, and
$vz_i=v(x_i+y_i)+b_{e_i}vt_0$, so pulling back carries the intercept
automatically and keys need no intercept correction. Such keys therefore already occur among
tables on equation tuples.

Take vertices to be key orbits under addition of a constant from $K$.
Choose the representative whose intercept has $K$-component zero. Restore
a table's output by adding its removed constant to the vertex label.
Every labeling then defines functions satisfying
\begin{equation}\label{eq:fold-share}
 F_Q(P+h e_0^\top)=F_Q(P)+h,\qquad
 F_U(P'\pi)=F_O(P').
\end{equation}
Here $e_0^\top$ is the first-coordinate functional.

\paragraph{The test and completeness.}
Sample $U\sim\omega^{\otimes k}$, a uniform table $P$, noise $a\sim\mu$, and a uniform
$l\in\F_2^m$, where $m=1+2k$. Require
\begin{equation}\label{eq:matrix-test}
 F_U(P)=F_U(P+a l^\top).
\end{equation}
Equality of restored outputs is a translation constraint on the two vertex
labels. The sampling probabilities are its rational weights. Call this
weighted translation game $\mathcal H$. Loops and repeated edges are allowed
for now.

Given a global assignment satisfying at least $1-\xi$ of the equations,
answer each retained variable by its global bit, and each retained equation
by the global first two bits and the required third bit. Label a key
representative by evaluating its table on this answer and applying $C$.
Equivariance makes restoration consistent. With probability at least
$1-k\xi$, every equation in $U$ is satisfied. The answers used for all keys of tables on $U$ are then
projections of one legal vector $z_U$. The compared values are
\[
 C(Pz_U),\qquad C(Pz_U+a(l^\top z_U)).
\]
Since $z_U\ne0$, the first argument is uniform and $l^\top z_U$ is an
independent fair bit. Consequently
\begin{equation}\label{eq:completeness}
 \val(\mathcal H)\ge1-k\xi-p_*/2.
\end{equation}

\paragraph{Why a gadget is needed.}
Without one ($\mathcal V=K$, $C$ the identity, $a$ uniform on $K$), the
labeling above is $F_U(P)=Pz_U$ on tuples whose equations are all
satisfied, and $F_U(P+al^\top)=Pz_U+a(l^\top z_U)$. As $l^\top z_U$ is a
fair bit independent of $a$, the test \eqref{eq:matrix-test} accepts there
with probability exactly $\frac12+\frac1{2|K|}$. The gadget's $C$ absorbs
the noise, giving \eqref{eq:completeness}, yet every linear family whose
restriction to $K$ has large dimension still detects it, which
Lemma~\ref{lem:advice} uses.

\section{The latent noise gadget}
\label{latent:section}

We prove the contract \eqref{eq:gadget-contract}, with $r_*$ depending only on $p_*$.

\subsection{One quadratic block}

Take $F=\F_q$, where $q=2^d$, and put $X=B_{\rm blk}=F^3$ and $d_{\rm blk}=3d$.
Identify their binary duals with $F^3$ using the nondegenerate pairing
$\Tr(u\cdot v)$, where $\Tr:F\to\F_2$ is the field trace. Define
\[
 Q(x)=(x_2x_3,x_1x_3,x_1x_2),\qquad
 C_{\rm blk}(x,y)=y+Q(x).
\]
For each one-dimensional $F$-subspace $A$ of $F^3$ fix $v\in A\setminus\{0\}$, and let
\[
 U_A=\{(tv,w):\phi_v(t,w)=0\},\qquad
 \phi_v(t,w)=v\cdot w+t^2v_1v_2v_3\quad(t\in F,\ w\in F^3).
\]
As $\phi_v$ is binary-linear and onto $F$, $U_A$ is a binary subspace of
dimension $4d-d=d_{\rm blk}$. One coordinate,
$(x_2+a_2)(x_3+a_3)+x_2x_3+a_2a_3=x_2a_3+x_3a_2$, gives by symmetry
$Q(x+a)+Q(x)+Q(a)=x\times a$, the cross product (signs do not matter in
characteristic two). Hence the update $(tv,w)$ changes the output by
$w+t^2Q(v)+t\,x\times v$, whose dot product with $v$ is $\phi_v(t,w)$, as
$v\cdot Q(v)=v_1v_2v_3$ and $v\cdot(x\times v)=0$. For a fixed input, each
$t$ thus has exactly one $w$ keeping the output, and it lies in $U_A$. So
exactly $q-1$ of the $q^3-1$ nonzero updates keep it, and a uniform nonzero
update changes it with probability $1-\theta$, where
$\theta=(q-1)/(q^3-1)=(q^2+q+1)^{-1}$ is one over the number of field
lines.

Fix a binary subspace $S\le B_{\rm blk}^*$ and a linear extension
$z\mapsto(\gamma_z,z)\in X^*\oplus B_{\rm blk}^*$, before choosing $A$,
with $\gamma_z(x)=\Tr(g(z)\cdot x)$ for binary-linear $g:S\to F^3$. For
$z\ne0$,
\begin{equation}
 (\gamma_z,z)|_{U_A}=0
 \quad\Longleftrightarrow\quad
 A=Fz\ \text{ and }\ (z\cdot g(z))^2=z_1z_2z_3.
 \label{latent:kernel}
\end{equation}
Indeed, the binary functionals vanishing on $\ker\phi_v$ are
$\Tr(c\,\phi_v)$, $c\in F$. Matching one with $(t,w)\mapsto\Tr(t\,g(z)\cdot
v+z\cdot w)$ forces $z=cv$, so $A=Fz$, and with $t=cs$ leaves
$\Tr(s\,g(z)\cdot z+s^2z_1z_2z_3)=0$ for all $s\in F$. Since
$\Tr(s\,g(z)\cdot z)=\Tr(s^2(g(z)\cdot z)^2)$ and squaring permutes $F$,
nondegeneracy of the trace gives \eqref{latent:kernel}.

\subsection{Generic spaces work for every extension}

Call an $r$-dimensional binary subspace $S\le F^3$ \emph{generic} if, for
every binary-linear $g:S\to F^3$, the alignment equation in
\eqref{latent:kernel} has at most $3r$ nonzero solutions.
For each fixed $r$, a uniform $r$-dimensional subspace is generic with
probability tending to one as $q$ tends to infinity.

Use independent uniform parameters $y_{ij}\in F$, and write
\[
 z_i(t)=\sum_{j=1}^r y_{ij}t_j,\qquad
 c_t=z_1(t)z_2(t)z_3(t),
 \qquad t\in\F_2^r.
\]
Conditional on injectivity of $t\mapsto z(t)$, which fails with
probability at most $2^rq^{-3}$, its image is a uniform $r$-space.

Fix a set $E\subseteq\F_2^r\setminus\{0\}$ with $|E|>3r$, and form
\[
 M_{t,(i,j)}=t_jz_i(t)^2\qquad(t\in E,\ 1\le i\le3,\ 1\le j\le r).
\]
Writing $g_i(z(t))=\sum_ja_{ij}t_j$, the squared alignment equation at $t$
is $\sum_{i,j}M_{t,(i,j)}a_{ij}^2=c_t$. So an extension aligning on $E$
gives $c=Mu$ with $u_{ij}=a_{ij}^2$; this single system handles every
observer.

Over the field $\mathbb L=\F_2(y_{ij})$, with the parameters as
indeterminates, $c$ is not in the column space $\mathcal C$ of $M$. Indeed,
suppose $c=Mu$ with $u\in\mathbb L^{3r}$. Differentiate formally in the parameters, extending partial
derivatives to $\mathbb L$ by the quotient rule. The entries of $M$
involve only squares $y_{ij}^2$, so in characteristic two every
derivative of $M$ vanishes. Hence, for all $j,k,l$,
\[
 (t_jt_kt_l)_{t\in E}
 =\partial_{y_{1j}}\partial_{y_{2k}}\partial_{y_{3l}}c
 =M\,\partial_{y_{1j}}\partial_{y_{2k}}\partial_{y_{3l}}u\in \mathcal C.
\]
The three differentiated variables are distinct even when $j=k=l$, so no
factorial vanishes; this case gives $t_j^3=t_j$, so every coordinate vector
$(t_j)_{t\in E}$ lies in $\mathcal C$. Multiplying a column of $M$ by a
coordinate stays in $\mathcal C$, since
$t_aM_{t,(i,j)}=\sum_hy_{ih}^2\,t_at_jt_h$ is a combination of the cubic
vectors above. Thus $\mathcal C$ is closed under coordinatewise
multiplication by each $t_a$, and so contains every positive-degree
monomial. The point indicator of a nonzero $e\in E$ vanishes at $0$, so it
has no constant term, and its values on $E$, the unit vector at $e$, lie in
$\mathcal C$; hence $\dim \mathcal C=|E|>3r$, impossible for $3r$ columns.

So some $\lambda$ over $\mathbb L$ has $\lambda M=0$ and
$\lambda c\ne0$; scale to polynomials. Parameters admitting $c=Mu$ give
$\lambda c=0$, as $\lambda M=0$ identically. By Schwartz--Zippel, the
nonzero polynomial $\lambda c$, fixed with $E$ and independent of $q$,
vanishes at uniform $F$-parameters with probability at most $\deg(\lambda
c)/q$. A union over the $E$ of size $3r+1$, whose number depends only on
$r$, and the conditioning on injectivity, a factor tending to one, finish
the proof.


\subsection{Recursion reduces the nonlinear error}

Each level has input $V_t=W_t\oplus B_{\rm blk}$ and output map
$C_t(w,b)=b+c_t(w)$. Start with $W_0=0$, $V_0=B_{\rm blk}$ and $C_0(b)=b$.
At each later height use one child for each pair in $\mathcal
I=\{(A,J):J:B_{\rm blk}\xrightarrow{\sim}U_A\}$ with binary isomorphisms,
and put
\[
 W_t=W_{t-1}^{\mathcal I}\oplus X,\qquad
 C_t((w_i)_i,x,b)=C_{\rm blk}\Bigl((x,b)+\sum_{i\in\mathcal I}J_ic_{t-1}(w_i)\Bigr),
\]
which has the required form with $c_t((w_i),x)=C_t((w_i),x,0)$, so
\begin{equation}
 C_t(w,b+h)=C_t(w,b)+h\qquad(h\in B_{\rm blk}).
 \label{latent:recursive-symmetry}
\end{equation}
For child $i$ define the linear injection $\iota_i:V_{t-1}\to V_t$,
$\iota_i(u,v)=\bigl((\ind_{j=i}u)_j,\,J_iv\bigr)$, with $J_iv\in U_A$
filling $(x,b)$. Let $\xi_0$ be uniform on $B_{\rm blk}$ and
$\xi_t=\iota_i\xi_{t-1}$ for uniform $i$ and independent $\xi_{t-1}$. Each
line has the same number of isomorphisms, so uniform $i$ is a uniform line
followed by a fresh uniform isomorphism.

Let $\Delta=v+c_{t-1}(w_i+u)+c_{t-1}(w_i)$ be the output change of the
chosen child under noise $(u,v)$; it does not involve the child's last
coordinate. For uniform input the child coordinates are uniform and
independent of $i$ and the noise, so $\Delta$ has $i$-independent law,
nonzero with probability $p_{t-1}$. Under $\xi_t$,
$\beta=(x,b)+\sum_jJ_jc_{t-1}(w_j)$ moves to $\beta+J_i\Delta$, with
$\beta$ uniform conditional on $i$, child noise and coordinates since
$(x,b)$ is uniform and independent of them. Given nonzero $\Delta$ and the
line, the fresh isomorphism makes $J_i\Delta$ uniform nonzero in $U_A$. The
block count gives $p_t=(1-\theta)p_{t-1}\le(1-\theta)^t$, where $p_t$ is
the probability that $C_t$ changes under its noise, on a uniform input
independent of the noise.

\subsection{Linear rank survives the recursion}

Put $H_0=0$ and $H_r=\sum_{i=1}^r1/i$. Let $G\le V_t^*$, for $t\ge1$, write
$(\gamma_\psi,z_\psi)$ for the restriction of $\psi\in G$ to $X\oplus
B_{\rm blk}$, and let $S=G|_{B_{\rm blk}}$ have dimension $r\ge1$. For
child $i=(A,J)$, $(\psi\circ\iota_i)(0,v)=(\gamma_\psi,z_\psi)(Jv)$, so
with $E_A=\{(\gamma_\psi,z_\psi)|_{U_A}:\psi\in G\}$,
$(\iota_i^*G)|_{B_{\rm blk}}=J^*E_A$. A linear section $\varsigma:S\to G$
of restriction gives the extension $z\mapsto(\gamma_{\varsigma(z)},z)$,
whose restrictions to $U_A$ lie in $E_A$. If $k_A$ is the kernel dimension
of this restriction map $S\to E_A$, then $\dim E_A\ge r-k_A$. By
\eqref{latent:kernel}, every nonzero kernel element lies on $A$ and is
aligned, and each aligned element lies on exactly one line. Let $N_S$ be
the largest number of aligned elements over all extensions; $N_S\le3r$ if
$S$ is generic, and $N_S<2^r$ always. Then $\sum_A(2^{k_A}-1)\le N_S$.
Since $1/(r-j)\le(j+1)/r\le2^j/r$ for $0\le j<r$, $H_r-H_{r-h}\le(2^h-1)/r$
for $0\le h\le r$; averaging over the $\theta^{-1}$ lines gives
\begin{equation}
 \E_AH_{\dim E_A}\ge H_r-\theta N_S/r.
 \label{latent:good-loss}
\end{equation}

Choose $r_0\ge1$ with $e^{1-H_{r_0}/8}\le p_*$, possible as $H_{r_0}\to\infty$,
and set $r_*=r_0$; then let the field degree $d$ be any integer large enough
for $2^{r_0}\theta\le\frac14$ and nongeneric proportion at most $2^{-r_0}$ at
each rank $1$ to $r_0$. Say $G$
detects $\xi$ if some element is nonzero at $\xi$. We claim that, for every
$t\ge0$ and every $G\le V_t^*$ with $\dim G\le r_0$, $S=G|_{B_{\rm blk}}$,
$r=\dim S$ ($N_S/r=0$ if $r=0$),
\begin{equation}
 2H_{r_0}\Pr[G\text{ detects }\xi_t]\ge H_r-\theta N_S/r-4\theta t.
 \label{latent:induction}
\end{equation}
If $r=0$ the right side is not positive; if $t=0$, $G=S$ detects uniform
$\xi_0$ with probability at least $\frac12$, and $H_{r_0}\ge H_r$ as $r\le
r_0$. Let $t,r\ge1$. Since $\psi(\xi_t)=(\psi\circ\iota_i)(\xi_{t-1})$, the
left side averages over uniform $i$ the left side for $\iota_i^*G$ at
height $t-1$; that family has dimension at most $r_0$ and restriction
$S'=J^*E_A$ of rank $r'=\dim E_A$. Given $A$, fresh $J$ makes $S'$ uniform
at its rank $r'\le r_0$, and $N_{S'}/r'<2^{r_0}$, so $\E[N_{S'}/r'\mid
A]\le3+2^{r_0}2^{-r_0}=4$. With \eqref{latent:good-loss},
\[
 \E_i\bigl[H_{r'}-\theta N_{S'}/r'\bigr]\ge H_r-\theta N_S/r-4\theta,
\]
so \eqref{latent:induction} at $t-1$ for each child gives it at $t$.

\paragraph{The output.} Put $\ell=3d$, $s=\lfloor
H_{r_0}/(8\theta)\rfloor\ge1$ (as $8\theta\le2^{1-r_0}\le1$), $\mathcal
V=V_s$, $W=W_s$, $C=C_s$, and $K=B_{\rm blk}=\F_2^\ell$ in a fixed binary
basis of $F^3$; let $\mu$ be the law of $\xi_s$. By
\eqref{latent:recursive-symmetry}, $C(x+h)=C(x)+h$ for $h\in K$, and
$p_s\le(1-\theta)^s\le e^{-\theta s}\le e^{\theta-H_{r_0}/8}\le p_*$. If
$\dim(G|_K)\ge r_0$, the span of $r_0$ elements of $G$ with independent
restrictions to $K$ satisfies \eqref{latent:induction} with $\theta
N_S/r_0<2^{r_0}\theta\le\frac14H_{r_0}$ and $4\theta s\le\frac12H_{r_0}$;
so it, and hence $G$, detects $\xi_s$ with probability at least $\frac18$.
All spaces, maps and noise choices are finite, and uniform choices and
their pushforwards have rational probabilities. This proves the contract
\eqref{eq:gadget-contract}.


\section{Soundness}\label{sec:decoders}

Fix $p_*$ and its gadget threshold $r_*$. Set $\eta=2^{-r_*}/16$.
Apply Input 2 at retention $\eta/4$ to obtain $\alpha,r_{\rm s}$. Choose
$\ell$ divisible by $3$ and large enough for that input and the gadget, and so that
\begin{equation}\label{eq:matrix-parameters}
 2^{-(\ell-r_{\rm s})}\le\alpha/4,\qquad 2^{-\ell}<\eta/4.
\end{equation}
Now fix the gadget and splitting $\mathcal V=K\oplus W$. Only $k$ grows
from now on. Write a table as $P=(M,T)$ with values in $K\oplus W$.

\paragraph{The precise experiment.}
Sample $U\sim\omega^{\otimes k}$. Independently in each position, retain
the equation with probability $1-\beta$, or a uniform one of its variables
with probability $\beta=k^{-2/3}$. This gives $O$ and $\pi:L_U\to L_O$.
Sample a uniform $A:K\to\F_2^{r_{\rm s}}$, then independent columns of
$T_O:L_O\to W$ and $M_O:L_O\to K$, all uniform except at singleton
coordinates: there the column of $T_O$ is zero and that of $M_O$ is
uniform in $\ker A$. Put $S_O=AM_O$. The messages are exactly
\begin{equation}\label{eq:messages}
 (U,A,T_O\pi,S_O\pi),\qquad (O,A,T_O,S_O).
\end{equation}
Neither player receives $M_O$ or $\pi$. Each may use the fixed function
family and independent private randomness, but only its own message.

\begin{proposition}[Soundness]\label{prop:sound}
There is $k_1$, depending only on the fixed auxiliary data, such that
$\val(\mathcal H)<99/100$ for every integer $k\ge k_1$ and every parity
instance with $\OPT(\mathcal E)\le6/7$, as in a NO instance with $\xi<1/100$.
\end{proposition}

\noindent Choose $k_0$ so that every $k\ge k_0$, with $m=1+2k$, gives
$2^{-m}\le\eta/4$ and makes $\ell+m$ exceed Input~2's ambient-dimension
threshold at retention $\eta/4$; the condition on $v_k$ is added below.
If $k\ge k_0$ and a labeling passes \eqref{eq:matrix-test} with probability
at least $99/100$, Sections~\ref{sec:noise}--\ref{sec:decoder} give two
strategies of this kind, returning legal answers (not necessarily
efficiently), with
\begin{equation}\label{eq:agreement}
 \Pr[\pi a_U=a_O]\ge\gamma,
\end{equation}
where $\gamma>0$, given in \eqref{eq:gamma}, depends only on the fixed
auxiliary data. If $\OPT(\mathcal E)\le6/7$, Section~\ref{sec:clean} bounds
the same probability by $\exp(-k^{1/3}/7056)$ for every such pair. Take
$k_1\ge k_0$ with $\exp(-k_1^{1/3}/7056)<\gamma$; for $k\ge k_1$ the two
bounds contradict each other, so no labeling passes with probability at
least $99/100$, that is, $\val(\mathcal H)<99/100$.

\subsection{Replacing the noise and extracting constant outputs}\label{sec:noise}

For a real function $f$ on a finite binary space, put
$\widehat f(\phi)=\E_x f(x)(-1)^{\phi(x)}$ for each functional $\phi$.
A nonzero character averages to zero: translation by a point where it is
one reverses its sign. Thus characters form an orthonormal basis, Fourier
inversion holds, and Parseval gives $\sum_\phi\widehat f(\phi)^2=\E f^2$.

For fixed $(U,T)$ write $f_{U,T}(M)=F_U(M,T)$ for $M\in\Hom(\F_2^m,K)$.
Call a nonempty row fiber $AM=S$, with $A:K\to\F_2^{r_{\rm s}}$, good if at
most $r_{\rm s}$ column conditions $Mq_j=t_j$ cut out a nonempty slice of it
on which one value of $f_{U,T}$ has density at least $\alpha$; then
$(U,A,T,S)$ is good row advice. Let
$\mathcal R_{U,T}$ be the union of all good row fibers, over all $A$ and $S$,
and write $\|\mathcal S\|$ for the fraction of all matrices in a set $\mathcal S$.

\begin{lemma}[Good row fibers]\label{lem:advice}
For $k\ge k_0$, if a labeling of $\mathcal H$ passes \eqref{eq:matrix-test} with
probability at least $99/100$, then $\E_{U,T}\|\mathcal R_{U,T}\|\ge\eta$.
\end{lemma}

\begin{proof}
\emph{Noise comparison.} Fix $U$. For a noise law $\nu$ on $\mathcal V$ let
$\mathrm{acc}_\nu=\Pr[F_U(P)=F_U(P+al^\top)]$, with $P$ and $l$ uniform and
$a\sim\nu$, and let $\nu_K$ be uniform on $K$. For $u\in K$ put
$f_u=\ind_{F_U=u}$, so that $\ind_{F_U(P)=F_U(P')}=\sum_uf_u(P)f_u(P')$.
Expand each $f_u$ in characters $g=(g_1,\ldots,g_m)\in(\mathcal V^*)^m$ of
the table space. Orthogonality over $P$ and averaging over $l$ give
\[
 \mathrm{acc}_\nu=\sum_gw(g)\lambda_\nu(g),
 \qquad w(g)=\sum_u\widehat{f_u}(g)^2,
 \qquad
 \lambda_\nu(g)=\Pr_{a\sim\nu}[g_i(a)=0\text{ for all }i]\ge0,
\]
and Parseval makes the weights sum to $\sum_u\E f_u=1$. Call $g$ low when
the restrictions $g_i|_K$ span dimension less than $r_*$, and let
$m_{\rm low}$ be the mass on low frequencies. For the others, the family
$\Span\{g_i\}$ has restriction to $K$ of dimension at least $r_*$, and
some element of it is nonzero at $a$ exactly when some $g_i$ is. So detection
gives $\lambda_\mu\le7/8$, whence $\mathrm{acc}_\mu\le m_{\rm low}+\frac78(1-m_{\rm low})$.
On the low frequencies $\lambda_{\nu_K}=2^{-\dim\Span(g_i|_K)}\ge2^{-r_*}$, so
$\mathrm{acc}_{\nu_K}\ge2^{-r_*}m_{\rm low}\ge2^{-r_*}(8\,\mathrm{acc}_\mu-7)$.
The noise $\nu_K$ leaves $T$ fixed, so $\mathrm{acc}_{\nu_K}=\E_T\,\mathrm{acc}(f_{U,T})$,
where $\mathrm{acc}(f)=\Pr[f(M)=f(M+al^\top)]$ for $M$, $a\in K$ and $l$ uniform.
The bound is linear in $\mathrm{acc}_\mu$, whose average over $U$ is at least
$99/100$, so $\E_{U,T}\,\mathrm{acc}(f_{U,T})\ge0.92\cdot2^{-r_*}=14.72\eta>3\eta$.

\emph{Covering.} It remains to show $\mathrm{acc}(f)\le\eta+2\|\mathcal R\|$
for each $(U,T)$, with $f=f_{U,T}$ and $\mathcal R=\mathcal R_{U,T}$, since
averaging then gives $\E_{U,T}\|\mathcal R_{U,T}\|\ge(3\eta-\eta)/2=\eta$.
The \emph{step retention} $\mathrm{ret}(\mathcal S)$ of a nonempty set
$\mathcal S$ is the probability that the uniform rank-one step
$M\mapsto M+al^\top$, zero step included, keeps a uniform member inside
it. For each value $u$ put $\mathcal S_u=f^{-1}(u)\setminus\mathcal R$, and
suppose that some nonempty $\mathcal S_u$ has $\mathrm{ret}(\mathcal S_u)\ge\eta$.
The choice of $\ell$ and $k_0$ lets Input~2 apply at retention $\eta/4$
and makes the step zero with probability
$2^{-\ell}+2^{-m}-2^{-\ell-m}\le\eta/2$, so conditioning on nonzero factors
leaves retention at least $\eta/2$. Identify $M$ with its graph
\[
 \Gamma_M=\{(x,M^\top x):x\in\F_2^\ell\},\qquad
 \dim(\Gamma_M\cap \Gamma_N)=\ell-\rank(M-N).
\]
These form the matrix chart in the Grassmann graph on
$\F_2^\ell\oplus\F_2^m$. A chart vertex has
$(2^\ell-1)(2^m-1)$ chart neighbors and
$(2^\ell-1)(2^{m+1}-2)$ neighbors in total. For the latter count, choose
a hyperplane in the vertex and a different one-dimensional extension in
the quotient. For the former, use the unique nonzero rank-one
factorization over $\F_2$. Exactly half the neighbors are in the chart,
so the lifted $\mathcal S_u$ has Grassmann retention at least $\eta/4$.
Input~2 gives it relative density at least $\alpha$ in an interval, hence also in
the interval's intersection with the chart, since the chart contains the whole lifted set.
A basis $(d_i,s_i)$ of $A_0$ and a basis $(t_j,q_j)$ of $B_0^\perp$, for the
dot product under which $\Gamma_M^\perp=\{(Mv,v)\}$, express this
intersection as the nonempty slice
\[
 \{M:d_i^\top M=s_i^\top,\ Mq_j=t_j\},
 \qquad i\le n_{\rm row},\ j\le n_c,\ n_{\rm row}+n_c\le r_{\rm s}.
\]
Its row conditions say $A_*M=S_*$, where $A_*$ has the rows $d_i^\top$
padded with zeros. As $\mathcal S_u\subseteq f^{-1}(u)$, this row fiber is
good and lies in $\mathcal R$; yet the slice meets $\mathcal S_u$, which is
disjoint from $\mathcal R$. Hence every nonempty $\mathcal S_u$ has step
retention below $\eta$.

Let $M'=M+al^\top$ be the step from a uniform $M$; both endpoints are
uniform. An accepted step with neither endpoint in $\mathcal R$ stays in
one set $\mathcal S_u$, so its probability is
$\sum_u\|\mathcal S_u\|\,\mathrm{ret}(\mathcal S_u)\le\eta$, the sum over
nonempty sets. A step with an endpoint in $\mathcal R$ has probability at
most $2\|\mathcal R\|$.
\end{proof}

\noindent For each matrix in $\mathcal R_{U,T}$ of Lemma~\ref{lem:advice}, at least one
of the $2^{\ell r_{\rm s}}$ possible row maps witnesses goodness. Hence,
under a uniform full table and independent uniform $A$,
\begin{equation}\label{eq:good-advice}
 \Pr[(U,A,T,AM)\text{ is good row advice}]
 \ge\eta\,2^{-\ell r_{\rm s}}.
\end{equation}
\subsection{The private decoders and their agreement}\label{sec:decoder}

Fix $U$ and $A$, and put $n=|\mathcal V|$. Compare the law
$\mathcal L_{\rm pad}$ of the padded full table $(M_O\pi,T_O\pi)$, with its
intercept and averaged over the projection choices, with the law
$\mathcal L_{\rm unif}$ of a uniform table on $L_U$. Blocks are independent,
and a retained equation has uniform slopes. A singleton block has some
slope law $Q_A$, and $Q_A\le n^2P_0$ for the uniform law $P_0$ on
$\mathcal V^2$, so
$\chi^2(Q_A\Vert P_0):=\E_{P_0}(Q_A/P_0-1)^2=\E_{Q_A}(Q_A/P_0)-1<n^2$.
Mixing with weight $\beta$ multiplies the centered ratio by $\beta$, and
$1+\chi^2$ multiplies over independent blocks. The intercept changes
nothing: pulling back adds $b_{e_i}$ times each third-bit singleton column
to it, and an independent uniform intercept stays uniform after adding
terms that depend on the slopes. As Cauchy--Schwarz bounds total
variation by half the square root of $\chi^2$,
\begin{equation}\label{eq:tv}
 \bigl\|\mathcal L_{\rm pad}-\mathcal L_{\rm unif}\bigr\|_{\rm TV}
 \le v_k:=\frac12\sqrt{\bigl(1+n^2\beta^2\bigr)^k-1}.
\end{equation}
Averaging over $A$, the same bound compares the joint law of $A$ and the
padded table with an independent uniform pair. It does \emph{not}
compare joint laws that include $\pi$.
Put $H=\ker A$. By \eqref{eq:matrix-parameters}, $|H|\ge2^{\ell-r_{\rm s}}\ge4/\alpha$,
so $1/(|H|-1)\le2/|H|\le\alpha/2$. Put $b_*:=\eta\,2^{-2\ell r_{\rm s}}$. As
$k\beta^2=k^{-1/3}\to0$, enlarge $k_0$ so that $v_k\le\alpha b_*/8$ for
every $k\ge k_0$.

\paragraph{The strategies.}
Select a certificate deterministically from each good advice, using a
fixed order and the function family. Put $Z=\Span\{q_j\}$ and
\begin{equation}\label{eq:answer-list}
 \mathcal A=Z\cap\{v:e_0^\top v=1\},\qquad |\mathcal A|\le2^{r_{\rm s}}.
\end{equation}
On its message $(U,A,T_O\pi,S_O\pi)$, the first player samples uniformly
from this list, and gives any legal answer when the list is empty or the
advice is not good. On its message $(O,A,T_O,S_O)$, the second player
samples a uniform guess $u\in K$, a uniform $h\in H\setminus\{0\}$ and a
uniform $B\in\Hom(L_O,H)$. It fixes $M_0$ with $AM_0=S_O$ and, with
$h\lambda:x\mapsto\lambda(x)h$, puts
\[
 f(N)=\ind_{F_O(M_0+N,T_O)=u}-\ind_{F_O(M_0+N,T_O)=u+h},\qquad
 \varphi(\lambda)=f(B+h\lambda)
\]
for $N\in\Hom(L_O,H)$ and $\lambda\in L_O^*$. Indexing characters of
$L_O^*$ by $z\in L_O$, it returns $a_O=z$ with probability
$\widehat\varphi(z)^2/\E\varphi^2$, a law by Parseval, or a fixed legal
answer if $\varphi=0$. Folding \eqref{eq:fold-share} gives
$f(N+he_0^\top)=-f(N)$; substituting $\lambda+e_0^\top$ then
gives $\widehat\varphi(z)=-(-1)^{e_0^\top z}\widehat\varphi(z)$, so $a_O$
is legal. Each player uses only its own message and private randomness;
$u_R$, $J$, $Z'$, $X$, $\tilde N$ and $Y$ below enter only the analysis.

\paragraph{Agreement.}
Expose $D=(U,O,\pi,A,T_O,S_O,h)$. Given $D$, both messages are fixed,
so the first player's list choice, $N=M_O-M_0$ and the decoder's
$(u,B,z)$ are independent; $N$ is uniform on $\Hom(L_O,H)$, since at
singleton coordinates too the row fiber is $\ker A$. The first player's full row advice
$R=(U,A,T_O\pi,S_O\pi)$ is a function of $D$, and so is its certificate.
For good $R$, let $u_R$ be the first value, in the fixed order, of density
at least $\alpha$ on its chosen column slice. Let $J$ mean that $R$ is good
and its column conditions $M_O\pi q_j=t_j$ hold. They read
$N\pi q_j=t_j-M_0\pi q_j$, so $J$ is a function of $D$ and $N|_{Z'}$, where
$Z'=\pi Z$. Put $X=\E[f(N)\mid D,N|_{Z'}]$ for $u=u_R$ if $R$ is good, and
$X=0$ otherwise.

If $R$ is not good, $X=0$ and the bound below is trivial. Otherwise, given $D$, the guess $u=u_R$ is
right with probability $2^{-\ell}$; assume it. Take $\lambda$ uniform on
$L_O^*$, independent of $D$ and $B$, and put $\tilde N=B+h\lambda$, so that
$\varphi(\lambda)=f(\tilde N)$ and $\tilde N|_{Z'}$ is a function of $B$ and
$\lambda|_{Z'}$. Averaging $\varphi$ over each set
$\lambda|_{Z'}=\text{const}$ keeps exactly its coefficients at $Z'$, so
Parseval and $\E\varphi^2\le1$ give, even if $\varphi=0$,
\[
 \Pr[a_O\in Z'\mid D,u,B]\ge\sum_{z\in Z'}\widehat\varphi(z)^2
 =\E[Y^2\mid D,B],
 \qquad Y:=\E[f(\tilde N)\mid D,B,\lambda|_{Z'}].
\]
By the tower property and conditional Jensen,
\[
 \E[Y^2\mid D]\ge\E\bigl[\E[f(\tilde N)\mid D,\tilde N|_{Z'}]^2\bigm|D\bigr]
 =\E[X^2\mid D],
\]
the equality because, given $D$, the uniform $B$ makes $\tilde N$ uniform on
$\Hom(L_O,H)$, like $N$, and $f$ is fixed by $D$. Averaging over $B$ gives
$\Pr[a_O\in Z',\ u=u_R\mid D]\ge2^{-\ell}\E[X^2\mid D]$.
On this event $a_O=\pi z'$ for some $z'\in Z$, and $z'\in\mathcal A$, as
$a_O$ is legal and $\pi$ preserves the first coordinate. The first
player's independent uniform choice from $\mathcal A$ hits $z'$ with
probability at least $2^{-r_{\rm s}}$, so
$\Pr[\pi a_U=a_O\mid D]\ge c\,\E[X^2\mid D]$ with $c:=2^{-r_{\rm s}-\ell}$.

It remains to bound $\E[\ind_JX]=\E[\ind_Jf(N)]$, with $u=u_R$. For the
full table $(M,T)=(M_O\pi,T_O\pi)$, sharing gives
$f(N)=\ind_{F_U(M,T)=u_R}-\ind_{F_U(M,T)=u_R+h}$, and $J$ says that the
advice $(U,A,T,AM)$ is good and $Mq_j=t_j$. So the score $\ind_Jf(N)$ is one formula,
bounded by one, in $U$, $A$, $h$ and the full table, not involving $\pi$;
as $h$ depends only on $A$, \eqref{eq:tv} applies to it. Under the uniform
law, given good advice $R$, $M$ is uniform on its row fiber, each nonempty
column slice imposes at most $\ell r_{\rm s}$ additional scalar equations on
it, and $h$ is independent of $M$; so the formula has conditional mean at
least $\alpha-1/(|H|-1)\ge\alpha/2$ given $R$ and $J$; and,
as the slice has probability at least $2^{-\ell r_{\rm s}}$ in its fiber,
\eqref{eq:good-advice} gives $\Pr_{\rm unif}(J)\ge b_*$. Hence
$\E[\ind_Jf(N)]\ge\alpha b_*/2-2v_k\ge\alpha b_*/4$, and Cauchy--Schwarz
gives
\begin{equation}\label{eq:gamma}
 \Pr[\pi a_U=a_O]
 \ge c\,\E X^2
 \ge c\bigl(\E[\ind_JX]\bigr)^2
 \ge\frac{c\alpha^2b_*^2}{16}
 =\frac{\alpha^2\eta^2}{16}\,2^{-4\ell r_{\rm s}-r_{\rm s}-\ell}=:\gamma.
\end{equation}
This proves \eqref{eq:agreement} for the original unconditioned experiment.


\subsection{Clean coordinates force disagreement}\label{sec:clean}

Let $\OPT(\mathcal E)\le6/7$, and take any two strategies that use only their
own messages \eqref{eq:messages} and private randomness and return legal
answers. By \eqref{eq:base-gap} the base game has value at most $20/21$. Call the
singleton positions clean, and let $I$ be their set. Given $A$, positions
are independent and the law of a position's columns depends only on whether
it is a singleton. So the experiment can be sampled in two stages: first
$A$, $I$, all $T_O,M_O$ columns and equations outside $I$; then for each
$i\in I$ an independent base round $e_i\sim\omega$ with a uniform variable.
At clean positions $T_O,S_O$ columns vanish, so the pulled-back table does
not depend on that position. Given the first stage, the first message is a
function of the clean equations and the second of the clean variables, with
one answer coordinate per position even when identifiers repeat. With the first
stage as a fixed parameter, sent to neither player, and keeping their clean
answers, the strategies thus play
the $|I|$-fold base game legally, full agreement winning every round.
Mixing over private randomness, Input~3 gives, also for $I=\varnothing$,
\[
 \Pr[\pi a_U=a_O\mid\text{first stage}]\le(1-1/7056)^{|I|}.
\]
Since $|I|\sim\mathrm{Bin}(k,\beta)$,
\begin{equation}\label{eq:clean-upper}
 \Pr[\pi a_U=a_O]\le(1-\beta/7056)^k\le\exp(-k^{1/3}/7056)\longrightarrow0.
\end{equation}
This bounds the same \emph{unconditioned} experiment as
\eqref{eq:agreement} and proves Proposition~\ref{prop:sound}.

\section{Removing weights, amplifying, and writing the output}\label{sec:finish}

\paragraph{An unweighted simple bipartite graph.} Weights of $\mathcal E$
are multiples of $1/(3N_0)$ and $\mu$ has fixed denominator $D_\mu$, so all
$\mathcal H$ probabilities are multiples of $1/N_{\mathcal H}$ with
$N_{\mathcal H}=(3N_0)^k|\mathcal V|^m2^mD_\mu$ polynomial (only $N_0$
grows). Replace a constraint of weight $w$ by $wN_{\mathcal H}$ four-edge
paths with three fresh internal vertices: three identity permutations then
the original constraint. For fixed endpoint labels, all four edges can be
satisfied exactly when the original constraint holds, and otherwise exactly
three; disjoint interiors let optima combine. Hence for $G_{\rm bip}$,
\begin{equation}\label{eq:subdivision}
 \val(G_{\rm bip})=1-\frac{1-\val(\mathcal H)}4.
\end{equation}
Place original vertices and path midpoints on one side and the other
internal vertices on the other; loops become four-cycles, parallel copies
get disjoint interiors, so the graph is simple and bipartite. Reversing an edge
preserves binary translations.

\paragraph{Parameter order.} Choose positive rationals $\eps_0\le\eps$,
$\delta_0\le\delta$. By Proposition~\ref{prop:sound} and
\eqref{eq:subdivision}, the NO gap after subdivision is at least $1/400$. Since
\eqref{eq:repetition} is alphabet-independent, first fix $t$ with
\begin{equation}\label{eq:final-t}
 \left(1-\frac1{2\,560\,000}\right)^t\le\delta_0,
 \qquad 2\,560\,000=16\cdot400^2.
\end{equation}
Next choose a positive rational $p_*\le\eps_0/t$; then $r_*$, $\eta$,
$\alpha,r_{\rm s}$, $\ell$ and the gadget as in Section~\ref{sec:decoders};
then $k=k_1$ of Proposition~\ref{prop:sound}; finally a positive rational
$\xi<\min\{1/100,\eps_0/(2kt)\}$. All are uniform over the instance and the function family,
hardcoded into one machine for fixed targets, so their existence suffices;
no uniform computability in targets is claimed.

\paragraph{The final game.} Output $G=G_{\rm bip}^{\otimes t}$. In the YES
case \eqref{eq:completeness} gives $\val(\mathcal H)\ge1-\mathrm{err}$ with
$\mathrm{err}=k\xi+p_*/2\le\eps_0/t$; a good labeling in every
coordinate and a union bound give
\[
 \val(G)\ge1-t\,\mathrm{err}/4\ge1-\eps_0/4\ge1-\eps.
\]
In the NO case, \eqref{eq:repetition} and \eqref{eq:final-t} give
$\val(G)\le(1-1/2\,560\,000)^t\le\delta_0\le\delta$. A pair of product
vertices specifies at most one edge per coordinate, so $G$ is simple. It is
bipartite, and independent uniform base edges give uniform product edges. Its
constraints are translations over the fixed alphabet $K^t=\F_2^{\ell t}$.
Keys need only fixed-dimensional linear algebra, even comparing every pair
of keys is polynomial, and $t$ is fixed, so $G$ is explicit in polynomial
time, proving Theorem~\ref{thm:main}.

{\small
\medskip\noindent Departures from the source proof~\cite{UGT}, credits for each
simplification, the review history and the Lean check are listed in the
companion note~\cite{Notes}, deposited with this paper.\par
}

{\small
\begin{thebibliography}{9}
\bibitem{UGT}
OpenAI. \emph{The Unique Games Theorem}. Preprint, 23 September 2026,
58 pages. Pinned source revision:
\href{https://github.com/openai/math/tree/fd4aeeb2ee4fc729c18d98444fed42fd0529eeeb/preprints/The-Unique-Games-Theorem-September-23-2026}{GitHub source and paper}.

\bibitem{Hastad}
Johan H\aa stad. \emph{Some Optimal Inapproximability Results}.
Journal of the ACM 48(4), 798--859, 2001.
Input 1 uses Theorem 5.4 and its proof in the
\href{https://people.kth.se/~johanh/optimalinap.pdf}{author manuscript}.

\bibitem{KMS}
Subhash Khot, Dor Minzer, and Muli Safra.
\emph{Pseudorandom Sets in Grassmann Graph Have Near-Perfect Expansion}.
Annals of Mathematics 198(1), 1--92, 2023.
Input 2 uses Theorem 1.12 in the
\href{https://eccc.weizmann.ac.il/report/2018/006/revision/2/download}{second 2018 revision}.

\bibitem{DS}
Irit Dinur and David Steurer. \emph{Analytical Approach to Parallel
Repetition}. Proceedings of STOC, 624--633, 2014.
Input 3 uses Corollary 1.2 in the
\href{https://www.dsteurer.org/paper/productgames.pdf}{10 June 2014 author version}.

\bibitem{Fortnow}
Lance Fortnow (prepared with Claude). \emph{The Unique Games Theorem: An
Expository Account of the Proof}. 7 October 2026,
\href{https://github.com/fortnow/AI-Papers}{github.com/fortnow/AI-Papers}.

\bibitem{Elboim}
Ariel Elboim. \emph{A Compressed Presentation of OpenAI's Proof of the
Unique Games Theorem}. Preprint, 2026,
\href{https://doi.org/10.5281/zenodo.23264902}{doi:10.5281/zenodo.23264902}.

\raggedright\bibitem{Notes}
Ariel Elboim. \emph{Companion Note to ``The Unique Games Theorem: A Short
Self-Contained Proof'': Attribution, Provenance and Checks}. 2026.
Deposited with this paper at
\href{https://doi.org/\depositdoi}{doi:\depositdoi}.
\end{thebibliography}
}
\end{document}
